\documentclass[12pt]{article}

% Computational Physics by M. Newman
% Setup file for exercises

% Packages
\usepackage{mathpazo}
\usepackage{amsmath}
\usepackage{graphicx}
\usepackage{calrsfs}
\usepackage{fancyvrb}

% Formatting
\setlength{\textwidth}{17.5cm}
\setlength{\textheight}{24.0cm}
\setlength{\evensidemargin}{2.5cm}
\setlength{\oddsidemargin}{-0.6cm}
\setlength{\topmargin}{-2.3cm}
\setlength{\textfloatsep}{0pt}
\setcounter{secnumdepth}{0}
\renewcommand{\baselinestretch}{1.06}

% Macros
\newcommand{\dd}{\mathrm{d}}
\newcommand{\ii}{\mathrm{i}}
\newcommand{\e}{\mathrm{e}}
\newcommand{\half}{\tfrac12}
\newcommand{\av}[1]{\langle#1\rangle}
\newcommand{\var}{\mathop\mathrm{var}}
\newcommand{\Ord}{\mathrm{O}}
\renewcommand{\Re}{\mathop\mathrm{Re}\nolimits}
\renewcommand{\Im}{\mathop\mathrm{Im}\nolimits}
\newcommand{\mat}{\mathbf}
\renewcommand{\vec}{\mathbf}
\newcommand{\defn}{\textit}
\newcommand{\exskip}{\nopagebreak\medskip\noindent}
\newcommand{\pmstrut}{\rule{0pt}{13pt}}

% Environments
\newcounter{chapter}
\setcounter{chapter}{2}
\newcounter{exercise}
\renewcommand{\theexercise}{\arabic{chapter}.\arabic{exercise}}
\newenvironment{exercises}{\vspace{2ex plus0.5ex minus0.5ex}
  \renewcommand{\labelenumi}{\alph{enumi})}}%
    {\par\vspace{1ex plus0.5ex minus0.5ex}}
\newcommand{\exercise}{\refstepcounter{exercise}%
  \par\vspace{4ex plus1ex minus1ex}%
  \noindent\textbf{Exercise \theexercise: }}

\DefineVerbatimEnvironment{code}{Verbatim}{xleftmargin=\parindent}

\setcounter{chapter}{7}

\begin{document}

\noindent {\LARGE\textsc{Computational Physics}}\par
\bigskip
\noindent {\large\textsc{Exercises for Chapter \arabic{chapter}}}\par
\noindent\hrulefill

%%%%%%%%%%%%%%%%%%%%%%%%%%%%%%%%%%%%%%%%%%%%%%%%%%%%%%%%%%%%%%%%%%%%%%
%%%                                                                %%%
%%%     COMPUTATIONAL PHYSICS, M. NEWMAN, CHAPTER 7, EXERCISES     %%%
%%%                                                                %%%
%%%%%%%%%%%%%%%%%%%%%%%%%%%%%%%%%%%%%%%%%%%%%%%%%%%%%%%%%%%%%%%%%%%%%%

\begin{exercises}

%%% Exercise 7.1 %%%

\exercise \textbf{Fourier transforms of simple functions}

\exskip Write Python programs to calculate the
coefficients in the discrete Fourier transforms of the following periodic
functions sampled at $N=1000$ evenly spaced points, and make plots of their
amplitudes similar to the plot shown in Fig.~7.4:
\begin{enumerate}\setlength{\itemsep}{0pt}\setlength{\parskip}{0pt}
\item A single cycle of a square-wave with amplitude~1
\item The sawtooth wave $y_n=n$
\item The modulated sine wave $y_n = \sin(\pi n/N) \sin(20\pi n/N)$
\end{enumerate}
If you wish you can use the Fourier transform function from the file
\verb|dft.py| as a starting point for your program.


%%% Exercise 7.2 %%%

\exercise \textbf{Detecting periodicity}

\exskip In the on-line resources there is a file called
\verb|sunspots.txt|, which contains the observed number of sunspots on the
Sun for each month since January 1749.  The file contains two columns of
numbers, the first representing the month and the second being the sunspot
number.
\begin{enumerate}\setlength{\itemsep}{0pt}
\item Write a program that reads the data in the file and makes a graph of
  sunspots as a function of time.  You should see that the number of
  sunspots has fluctuated on a regular cycle for as long as observations
  have been recorded.  Make an estimate of the length of the cycle in
  months.
\item Modify your program to calculate the Fourier transform of the sunspot
  data and then make a graph of the magnitude squared~$|c_k|^2$ of the
  Fourier coefficients as a function of~$k$---also called the
  \defn{power spectrum} of the sunspot signal.  You should see that
  there is a noticeable peak in the power spectrum at a nonzero value
  of~$k$.  The appearance of this peak tells us that there is one frequency
  in the Fourier series that has a higher amplitude than the others around
  it---meaning that there is a large sine-wave term with this frequency,
  which corresponds to the periodic wave you can see in the original data.
\item Find the approximate value of $k$ to which the peak corresponds.
  What is the period of the sine wave with this value of~$k$?  You should
  find that the period corresponds roughly to the length of the cycle that
  you estimated in part~(a).
\end{enumerate}
This kind of Fourier analysis is a sensitive method for detecting
periodicity in signals.  Even in cases where it is not clear to the eye
that there is a periodic component to a signal, it may still be possible to
find one using a Fourier transform.


%%% Exercise 7.3 %%%

\exercise \textbf{Fourier transforms of musical instruments}

\exskip In the on-line resources you will
find files called \texttt{piano.txt} and \texttt{trumpet.txt}, which
contain data representing the waveform of a single note, played on,
respectively, a piano and a trumpet.
\begin{enumerate}\setlength{\itemsep}{0pt}
\item Write a program that loads a waveform from one of these files, plots
  it, then calculates its discrete Fourier transform and plots the
  magnitudes of the first $10\,000$ coefficients in a manner similar to
  Fig.~7.4.  Note that you will have to use a fast Fourier transform for
  the calculation because there are too many samples in the files to do the
  transforms the slow way in any reasonable amount of time.

  Apply your program to the piano and trumpet waveforms and discuss briefly
  what one can conclude about the sound of the piano and trumpet from the
  plots of Fourier coefficients.
\item Both waveforms were recorded at the industry-standard rate of
  $44\,100$ samples per second and both instruments were playing the same
  musical note when the recordings were made.  From your Fourier transform
  results calculate what note they were playing.  (Hint: The musical note
  middle~C has a frequency of $261\,$Hz.)
\end{enumerate}


%%% Exercise 7.4 %%%

\exercise \textbf{Fourier filtering and smoothing}

\exskip In the on-line resources you'll find a file called \verb|dow.txt|.
It contains the daily closing value for each business day from late 2006
until the end of 2010 of the Dow Jones Industrial Average, which is a
measure of average prices on the US stock market.

Write a program to do the following:
\begin{enumerate}\setlength{\itemsep}{0pt}
\item Read in the data from \verb|dow.txt| and plot them on a graph.
\item Calculate the coefficients of the discrete Fourier transform of the
  data using the function \verb|rfft| from \verb|numpy.fft|, which produces
  an array of $\half N+1$ complex numbers.
\item Now set all but the first 10\% of the elements of this array to zero
  (i.e.,~set the last 90\% to zero but keep the values of the first 10\%).
\item Calculate the inverse Fourier transform of the resulting array, zeros
  and all, using the function \verb|irfft|, and plot it on the same graph
  as the original data.  You may need to vary the colors of the two curves
  to make sure they both show up on the graph.  Comment on what you see.
  What is happening when you set the Fourier coefficients to zero?
\item Modify your program so that it sets all but the first 2\% of the
  coefficients to zero and run it again.
\end{enumerate}


%%% Exercise 7.5 %%%

\exercise If you have not done Exercise~7.4 you should do it before this
one.

\exskip The function $f(t)$ represents a square-wave with amplitude~1 and
frequency $1\,$Hz:
\begin{equation}
f(t) = \left\lbrace\begin{array}{rl}
          1  & \qquad\mbox{if $\lfloor 2t \rfloor$ is even,} \\
         -1  & \qquad\mbox{if $\lfloor 2t \rfloor$ is odd,}
       \end{array}\right.
\end{equation}
where $\lfloor x\rfloor$ means $x$ rounded down to the next lowest integer.
Let us attempt to smooth this function using a Fourier transform, as we did
in the previous exercise.  Write a program that creates an array of
$N=1000$ elements containing a thousand equally spaced samples from a
single cycle of this square-wave.  Calculate the discrete Fourier transform
of the array.  Now set all but the first ten Fourier coefficients to zero,
then invert the Fourier transform again to recover the smoothed signal.
Make a plot of the result and on the same axes show the original
square-wave as well.  You should find that the signal is not simply
smoothed---there are artifacts, wiggles, in the results.  Explain
briefly where these come from.

Artifacts similar to these arise when Fourier coefficients are discarded in
audio and visual compression schemes like those described in Section~7.3.1
and are the primary source of imperfections in digitally compressed sound
and images.


%%% Exercise 7.6 %%%

\exercise \textbf{Comparison of the DFT and DCT}

\exskip This exercise will be easier if you have already done Exercise~7.4.

Exercise~7.4 looked at data representing the variation of the Dow Jones
Industrial Average, colloquially called ``the Dow,'' over time.  The
particular time period studied in that exercise was special in one sense:
the value of the Dow at the end of the period was almost the same as at the
start, so the function was, roughly speaking, periodic.  In the on-line
resources there is another file called \verb|dow2.txt|, which also contains
data on the Dow but for a different time period, from 2004 until 2008.
Over this period the value changed considerably from a starting level
around 9000 to a final level around 14000.
\begin{enumerate}\setlength{\itemsep}{0pt}
\item Write a program similar to the one for Exercise~7.4, part~(e), in
  which you read the data in the file \verb|dow2.txt| and plot it on a
  graph.  Then smooth the data by calculating its Fourier transform,
  setting all but the first 2\% of the coefficients to zero, and inverting
  the transform again, plotting the result on the same graph as the
  original data.  As in Exercise~7.4 you should see that the data are
  smoothed, but now there will be an additional artifact.  At the beginning
  and end of the plot you should see large deviations away from the true
  smoothed function.  These occur because the function is required to be
  periodic---its last value must be the same as its first---so it needs to
  deviate substantially from the correct value to make the two ends of the
  function meet.  In some situations (including this one) this behavior is
  unsatisfactory.  If we want to use the Fourier transform for smoothing,
  we would certainly prefer that it not introduce artifacts of this kind.
\item Modify your program to repeat the same analysis using discrete cosine
  transforms.  You can use the functions from \verb|dcst.py| to perform the
  transforms if you wish.  Again discard all but the first 2\% of the
  coefficients, invert the transform, and plot the result.  You should see
  a significant improvement, with less distortion of the function at the
  ends of the interval.  This occurs because, as discussed at the end of
  Section~7.3, the cosine transform does not force the value of the
  function to be the same at both ends.
\end{enumerate}
It is because of the artifacts introduced by the strict periodicity
of the DFT that the cosine transform is favored for many technological
applications, such as audio compression.  The artifacts can degrade the
sound quality of compressed audio and the cosine transform generally gives
better results.

The cosine transform is not wholly free of artifacts itself however.  It's
true it does not force the function to be periodic, but it does force the
gradient to be zero at the ends of the interval (which the ordinary Fourier
transform does not).  You may be able to see this in your calculations for
part~(b) above.  Look closely at the smoothed function and you should see
that its slope is flat at the beginning and end of the interval.  The
distortion of the function introduced is less than the distortion in
part~(a), but it's there all the same.  To reduce this effect, audio
compression schemes often use overlapped cosine transforms, in which
transforms are performed on overlapping blocks of samples, so that the
portions at the ends of blocks, where the worst artifacts lie, need not be
used.


%%% Exercise 7.7 %%%

\exercise \textbf{Fast Fourier transform}

\exskip Write your own program to compute the fast Fourier transform for
the case where $N$ is a power of two, based on the formulas given in
Section~7.4.1.  As a test of your program, use it to calculate the Fourier
transform of the data in the file \verb|pitch.txt|, which can be found in
the on-line resources.  A plot of the data is shown in Fig.~7.3.  You
should be able to duplicate the results for the Fourier transform shown in
Fig.~7.4.

This exercise is quite tricky.  You have to calculate the
coefficients~$E^{(m,j)}_k$ from Eq.~(7.43) for all levels~$m$, which means
that first you will have to plan how the coefficients will be stored.
Since, as we have seen, there are exactly $N$ of them at every level, one
way to do it would be to create a two-dimensional complex array of size
$N\times(1+\log_2 N)$, so that it has $N$ complex numbers for each level
from zero to~$\log_2 N$.  Then within level~$m$ you have $2^m$ individual
transforms denoted by $j=0\ldots 2^m-1$, each with $N/2^m$ coefficients
indexed by~$k$.  A simple way to arrange the coefficients would be to put
all the $k=0$ coefficients in a block one after another, then all the $k=1$
coefficients, and so forth.  Then $E^{(m,j)}_k$ would be stored in the
$j+2^mk$ element of the array.

This method has the advantage of being quite simple to program, but the
disadvantage of using up a lot of memory space.  The array contains
$N\log_2N$ complex numbers, and a complex number typically takes sixteen
bytes of memory to store.  So if you had to do a large Fourier transform
of, say, $N=10^8$~numbers, it would take $16N\log_2N\simeq42$~gigabytes of
memory, which is much more than most computers have.

An alternative approach is to notice that we do not really need to store
all of the coefficients.  At any one point in the calculation we only need
the coefficients at the current level and the previous level (from which
the current level is calculated).  If one is clever one can write a program
that uses only two arrays, one for the current level and one for the
previous level, each consisting of $N$ complex numbers.  Then our transform
of $10^8$ numbers would require less than four gigabytes, which is fine on
most computers.

(There is a third way of storing the coefficients that is even more
efficient.  If you store the coefficients in the correct order, then you
can arrange things so that every time you compute a coefficient for the
next level, it gets stored in the same place as the old coefficient from
the previous level from which it was calculated, and which you no longer
need.  With this way of doing things you only need one array of $N$ complex
numbers---we say the transform is done ``in place.''  Unfortunately,
this in-place Fourier transform is much harder to work out and harder to
program.  If you are feeling particularly ambitious you might want to give
it a try, but it's not for the faint-hearted.)


%%% Exercise 7.8 %%%

\exercise \textbf{Diffraction gratings}

\exskip Exercise~5.19 (page~206) looked at the physics of diffraction
gratings, calculating the intensity of the diffraction patterns they
produce from the equation
\begin{displaymath}
I(x) = \biggl| \int_{-w/2}^{w/2} \sqrt{q(u)}\>
               \e^{\ii 2\pi xu/\lambda f} \>\dd u \biggr|^2,
\end{displaymath}
where $w$ is the width of the grating, $\lambda$~is the wavelength of the
light, $f$~is the focal length of the lens used to focus the image, and
$q(u)$ is the intensity transmission function of the diffraction grating at
a distance~$u$ from the central axis, i.e.,~the fraction of the incident
light that the grating lets through.  In Exercise~5.19 we evaluated this
expression directly using standard methods for performing integrals, but a
more efficient way to do the calculation is to note that the integral is in
fact just a Fourier transform.  Approximating the integral, as we did in
Eq.~(7.13), using the trapezoidal rule, with~$N$ points $u_n=nw/N-w/2$, we
get
\begin{align*}
\int_{-w/2}^{w/2} \sqrt{q(u)}\>\e^{\ii 2\pi xu/\lambda f} \>\dd u
  &\simeq {w\over N} \e^{\ii\pi wx/\lambda f}
     \sum_{n=0}^{N-1} \sqrt{q(u_n)} \>\e^{\ii 2\pi wxn/\lambda f N} \\
  &= {w\over N} \e^{\ii\pi k}
     \sum_{n=0}^{N-1} y_n \>\e^{\ii 2\pi kn/N},
\end{align*}
where $k=wx/\lambda f$ and $y_n=\sqrt{q(u_n)}$.  Comparing with Eq.~(7.15),
we see that the sum in this expression is equal to the complex
conjugate~$c_k^*$ of the $k$th coefficient of the DFT of~$y_n$.
Substituting into the expression for the intensity~$I(x)$, we then have
\begin{displaymath}
I(x_k) = {w^2\over N^2} |c_k|^2,
\end{displaymath}
where
\begin{displaymath}
x_k = {\lambda f\over w}\,k.
\end{displaymath}
Thus we can calculate the intensity of the diffraction pattern at the
points~$x_k$ by performing a Fourier transform.

There is a catch, however.  Given that $k$ is an integer, $k=0\ldots N-1$,
the points~$x_k$ at which the intensity is evaluated have spacing~$\lambda
f/w$ on the screen.  This spacing can be large in some cases, giving us
only a rather coarse picture of the diffraction pattern.  For instance, in
Exercise~5.19 we had $\lambda=500\,$nm, $f=1\,$m, and $w=200\,\mu$m, and
the screen was $10\,$cm wide, which means that $\lambda f/w = 2.5\,$mm and
we have only forty points on the screen.  This is not enough to make a
usable plot of the diffraction pattern.

One way to fix this problem is to increase the width of the grating from
the given value~$w$ to a larger value~$W>w$, which makes the spacing
$\lambda f/W$ of the points on the screen closer.  We can add the extra
width on one or the other side of the grating, or both, as we prefer,
but---and this is crucial---the extra portion added must be opaque, it must
not transmit light, so that the physics of the system does not change.  In
other words, we need to ``pad out'' the data points~$y_n$ that measure the
transmission profile of the grating with additional zeros so as to make the
grating wider while keeping its transmission properties the same.  For
example, to increase the width to~$W=10w$, we would increase the number~$N$
of points~$y_n$ by a factor of ten, with the extra points set to zero.  The
extra points can be at the beginning, at the end, or split between the
two---it will make no difference to the answer.  Then the intensity is
given by
\begin{displaymath}
I(x_k) = {W^2\over N^2} |c_k|^2,
\end{displaymath}
where
\begin{displaymath}
x_k = {\lambda f\over W}\,k.
\end{displaymath}

Write a Python program that uses a fast Fourier transform to calculate the
diffraction pattern for a grating with transmission function $q(u) = \sin^2
\alpha u$ (the same as in Exercise~5.19), with slits of width $20\,\mu$m
[meaning that $\alpha=\pi/(20\,\mu$m)] and parameters as above:
$w=200\,\mu$m, $W=10w=2\,$mm, incident light of wavelength
$\lambda=500\,$nm, a lens with focal length of 1~meter, and a screen
$10\,$cm wide.  Choose a suitable number of points to give a good
approximation to the grating transmission function and then make a graph of
the diffraction intensity on the screen as a function of position~$x$ in
the range $-5\,$cm${}\le x\le 5\,$cm.  If you previously did Exercise~5.19,
check to make sure your answers to the two exercises agree.


%%% Exercise 7.9 %%%

\exercise \textbf{Image deconvolution}

\exskip You've probably seen it on TV, in one of those crime drama shows.
They have a blurry photo of a crime scene and they click a few buttons on
the computer and magically the photo becomes sharp and clear, so you can
make out someone's face, or some lettering on a sign.  Surely (like almost
everything else on such TV shows) this is just science fiction?  Actually,
no.  It's not.  It's real and in this exercise you'll write a program that
does it.

When a photo is blurred each point on the photo gets smeared out
according to some ``smearing distribution,'' which is technically called a
\defn{point spread function}.  We can represent this smearing
mathematically as follows.  For simplicity let's assume we're working with
a black and white photograph, so that the picture can be represented by a
single function~$a(x,y)$ which tells you the brightness at each
point~$(x,y)$.  And let us denote the point spread function by~$f(x,y)$.
This means that a single bright dot at the origin ends up appearing as
$f(x,y)$ instead.  If $f(x,y)$ is a broad function then the picture is
badly blurred.  If it is a narrow peak then the picture is relatively
sharp.

In general the brightness~$b(x,y)$ of the blurred photo at point $(x,y)$ is
given by
\begin{displaymath}
b(x,y) = \int_0^K \int_0^L a(x',y') f(x-x',y-y') \>\dd x'\>\dd y',
\end{displaymath}
where $K\times L$ is the dimension of the picture.  This equation is called
the \defn{convolution} of the picture with the point spread
function.

Working with two-dimensional functions can get complicated, so to get the
idea of how the math works, let's switch temporarily to a one-dimensional
equivalent of our problem.  Once we work out the details in 1D we'll return
to the 2D version.  The one-dimensional version of the convolution above
would be
\begin{displaymath}
b(x) = \int_0^L a(x') f(x-x') \>\dd x'.
\end{displaymath}
The function $b(x)$ can be represented by a Fourier series as in Eq.~(7.5):
\begin{displaymath}
b(x) = \sum_{k=-\infty}^\infty
         \tilde{b}_k \exp\biggl( \ii {2\pi k x\over L} \biggr),
\end{displaymath}
where
\begin{displaymath}
\tilde{b}_k = {1\over L} \int_0^L b(x)
                         \exp\biggl( -\ii {2\pi k x\over L} \biggr) \>\dd x
\end{displaymath}
are the Fourier coefficients.  Substituting for $b(x)$ in this equation
gives
\begin{align*}
\tilde{b}_k &= {1\over L} \int_0^L \int_0^L a(x') f(x-x')
                         \exp\biggl( -\ii {2\pi k x\over L} \biggr)
                         \>\dd x'\>\dd x \\
            &= {1\over L} \int_0^L \int_0^L a(x') f(x-x')
                         \exp\biggl( -\ii {2\pi k (x-x')\over L} \biggr)
                         \exp\biggl( -\ii {2\pi k x'\over L} \biggr)
                         \>\dd x'\>\dd x.
\end{align*}
Now let us change variables to $X=x-x'$, and we get
\begin{displaymath}
\tilde{b}_k = {1\over L} \int_0^L a(x')
                         \exp\biggl( -\ii {2\pi k x'\over L} \biggr)
                         \int_{-x'}^{L-x'} f(X)
                         \exp\biggl( -\ii {2\pi k X\over L} \biggr) \>\dd X
                         \>\dd x'.
\end{displaymath}
If we make $f(x)$ a periodic function in the standard fashion by repeating
it infinitely many times to the left and right of the interval from 0
to~$L$, then the second integral above can be written as
\begin{align*}
\int_{-x'}^{L-x'} f(X) \exp\biggl( -\ii {2\pi k X\over L} \biggr) \>\dd X
&= \int_{-x'}^0 f(X) \exp\biggl( -\ii {2\pi k X\over L} \biggr) \>\dd X
   \\
&\hspace{5em}{} + \int_0^{L-x'} f(X) \exp\biggl( -\ii {2\pi k X\over L}
   \biggr) \>\dd X \\
&\hspace{-12em} {} = \exp\biggl( \ii {2\pi k L\over L} \biggr)
   \int_{L-x'}^L f(X) \exp\biggl( -\ii {2\pi k X\over L} \biggr) \>\dd X
   + \int_0^{L-x'} f(X) \exp\biggl( -\ii {2\pi k X\over L} \biggr) \>\dd X
   \\
&\hspace{-12em} {} = \int_0^L f(X)
                    \exp\biggl( -\ii {2\pi k X\over L} \biggr) \>\dd X,
\end{align*}
which is simply $L$ times the Fourier transform~$\tilde{f}_k$ of~$f(x)$.
Substituting this result back into our equation for~$\tilde{b}_k$ we then
get
\begin{displaymath}
\tilde{b}_k = \int_0^L a(x')
              \exp\biggl( -\ii {2\pi k x'\over L} \biggr)
              \tilde{f}_k \>\dd x'
            = L\,\tilde{a}_k\tilde{f}_k.
\end{displaymath}
In other words, apart from the factor of~$L$, the Fourier transform of the
blurred photo is the product of the Fourier transforms of the unblurred
photo and the point spread function.

Now it is clear how we deblur our picture.  We take the blurred
picture and Fourier transform it to get $\tilde{b}_k =
L\,\tilde{a}_k\tilde{f}_k$.  We also take the point spread function and
Fourier transform it to get~$\tilde{f}_k$.  Then we divide one by the
other:
\begin{displaymath}
{\tilde{b}_k\over L\tilde{f}_k} = \tilde{a}_k
\end{displaymath}
which gives us the Fourier transform of the \emph{unblurred} picture.
Then, finally, we do an inverse Fourier transform on $\tilde{a}_k$ to get
back the unblurred picture.  This process of recovering the unblurred
picture from the blurred one, of reversing the convolution process, is
called \defn{deconvolution}.

Real pictures are two-dimensional, but the mathematics follows through
exactly the same.  For a picture of dimensions $K\times L$ we find that the
two-dimensional Fourier transforms are related by
\begin{displaymath}
\tilde{b}_{kl} =  KL\tilde{a}_{kl}\tilde{f}_{kl}\,,
\end{displaymath}
and again we just divide the blurred Fourier transform by the Fourier
transform of the point spread function to get the Fourier transform of the
unblurred picture.

In the digital realm of computers, pictures are not pure functions~$f(x,y)$
but rather grids of samples, and our Fourier transforms are discrete
transforms not continuous ones.  But the math works out the same again.

The main complication with deblurring in practice is that we don't usually
know the point spread function.  Typically we have to experiment with
different ones until we find something that works.  For many cameras it's a
reasonable approximation to assume the point spread function is Gaussian:
\begin{displaymath}
f(x,y) = \exp\biggl( -{x^2+y^2\over2\sigma^2} \biggr),
\end{displaymath}
where $\sigma$ is the width of the Gaussian.  Even with this assumption,
however, we still don't know the value of~$\sigma$ and we may have to
experiment to find a value that works well.  In the following exercise, for
simplicity, we'll assume we know the value of~$\sigma$.

\begin{enumerate}\setlength{\itemsep}{0pt}
\item On the web site you will find a file called \verb|blur.txt| that
  contains a grid of values representing brightness on a black-and-white
  photo---a badly out-of-focus one that has been deliberately blurred using
  a Gaussian point spread function of width $\sigma=25$.  Write a program
  that reads the grid of values into a two-dimensional array of real
  numbers and then draws the values on the screen of the computer as a
  density plot.  You should see the photo appear.  If you get something
  wrong it might be upside-down.  Work with the details of your program
  until you get it appearing correctly.  (Hint: The picture has the sky,
  which is bright, at the top and the ground, which is dark, at the
  bottom.)
\item Write another program that creates an array, of the same size as the
  photo, containing a grid of samples from drawn from the Gaussian~$f(x,y)$
  above with $\sigma=25$.  Make a density plot of these values on the
  screen too, so that you get a visualization of your point spread
  function.  Remember that the point spread function is periodic (along
  both axes), which means that the values for negative $x$ and $y$ are
  repeated at the end of the interval.  Since the Gaussian is centered on
  the origin, this means there should be bright patches in each of the four
  corners of your picture, something like this:
 \medskip
\begin{center}
\includegraphics[width=5cm,clip=true]{psf.eps}
\end{center}
\smallskip

\item Combine your two programs and add Fourier transforms using the
  functions \verb|rfft2| and \verb|irfft2| from \verb|numpy.fft|, to make a
  program that does the following:
  {\renewcommand{\labelenumii}{\roman{enumii})\ }
\begin{enumerate}\setlength{\itemsep}{0pt}\setlength{\parskip}{0pt}
\item Reads in the blurred photo
\item Calculates the point spread function
\item Fourier transforms both
\item Divides one by the other
\item Performs an inverse transform to get the unblurred photo
\item Displays the unblurred photo on the screen
\end{enumerate}}
\noindent When you are done, you should be able to make out the scene in
the photo, although probably it will still not be perfectly sharp.

Hint: One thing you'll need to deal with is what happens when the Fourier
transform of the point spread function is zero, or close to zero.  In that
case if you divide by it you'll get an error (because you can't divide by
zero) or just a very large number (because you're dividing by something
small).  A workable compromise is that if a value in the Fourier transform
of the point spread function is smaller than a certain amount~$\epsilon$
you don't divide by it---just leave that coefficient alone.  The value of
$\epsilon$ is not very critical but a reasonable value seems to be $10^{-3}$.
\item Bearing in mind this last point about zeros in the Fourier transform,
  what is it that limits our ability to deblur a photo?  Why can we not
  perfectly unblur any photo and make it completely sharp?
\end{enumerate}

We have seen this process in action here for a normal snapshot, but it is
also used in many physics applications where one takes photos.  For
instance, it is used in astronomy to enhance photos taken by telescopes.
It was famously used with images from the Hubble Space Telescope
after it was realized that the telescope's main mirror had a serious
manufacturing flaw and was returning blurry photos---scientists managed to
partially correct the blurring using Fourier transform techniques.

\end{exercises}

\end{document}
